\section*{Chapter \ref{ch:b3:statistical-thermo-applied} --- Statistical Thermodynamics Applied}

\begin{solution}{exo:b3:statistical-thermo-applied:1}
$\frac32R + R + (3 \times 3 - 5)R = 6.5R = \qty{54.0}{J.K^{-1}.mol^{-1}}$. At room temperature translation
and rotation are classical, the two stretches are frozen and the doubly
degenerate bend, the lowest vibration, is partly excited.
\end{solution}

\begin{solution}{exo:b3:statistical-thermo-applied:2}
$x = 1$: $\eu/(\eu - 1)^2 = 0.921$. $x = 10$: $100\eu^{10}/(\eu^{10} - 1)^2 \approx 100\eu^{-10} = \num{4.5e-3}$.
\end{solution}

\begin{solution}{exo:b3:statistical-thermo-applied:3}
\qty{20}{K}: $x_{\mathrm{para}} = 1/(1 + 3 \times 3\eu^{-2 \times 85.35/20}) = 1/(1 + 9 \times \num{1.96e-4}) = 0.998$.
\qty{77}{K}: even sum $1 + 5\eu^{-6.651} = 1.0065$; odd sum $3(3\eu^{-2.217} + 7\eu^{-13.30}) = 0.9806$;
$x_{\mathrm{para}} = 1.0065/1.9871 = 0.51$.
\end{solution}

\begin{solution}{exo:b3:statistical-thermo-applied:4}
For $I = 1$ there are $3 \times 2 = 6$ symmetric and $3 \times 1 = 3$ antisymmetric spin states;
deuterons are bosons, so the total wavefunction is symmetric: symmetric spin
states with even $J$ (ortho, weight 6), antisymmetric with odd $J$ (para, weight
3). Ortho-deuterium, which contains $J = 0$, is stable at 0 K; at room temperature
ortho:para = 2:1.
\end{solution}

\begin{solution}{exo:b3:statistical-thermo-applied:5}
$\Delta_rE_0 = 2 \times 1883.7 - 2170.3 - 1542.3 = \qty{54.8}{cm^{-1}}$; $\eu^{-54.8/207.2} = 0.768$, and
$4 \times 0.768 = 3.07$. The full value, 3.26, is 6\,\% higher: the translational, rotational and
vibrational factors only cancel exactly in the classical limit, and the
rotation of these light molecules is still partly quantised at \qty{298}{K}.
\end{solution}

\begin{solution}{exo:b3:statistical-thermo-applied:6}
\ce{N2O}: two orientations, $R\ln2 = \qty{5.76}{J.K^{-1}.mol^{-1}}$. \ce{CH3D}: four, $R\ln4 =
\qty{11.53}{J.K^{-1}.mol^{-1}}$ (upper limits, reached if the arrangements are random).
\end{solution}

\begin{solution}{exo:b3:statistical-thermo-applied:7}
$x = 797.6/300 = 2.659$: Einstein function $x^2\eu^{-x}/(1 - \eu^{-x})^2 = 0.572$; $C_V/R = 3.072$,
$C_p = 4.072R = \qty{33.86}{J.K^{-1}.mol^{-1}}$, 0.4\,\% below the table (anharmonicity).
\end{solution}

\begin{solution}{exo:b3:statistical-thermo-applied:8}
$x = 1.029$: $x^2\eu^{-x}/(1 - \eu^{-x})^2 = 0.916$, i.e. \qty{7.62}{J.K^{-1}.mol^{-1}}, 92\,\% of $R$: the
vibration of iodine is almost classical at room temperature.
\end{solution}

\begin{solution}{exo:b3:statistical-thermo-applied:9}
B is richer in \ce{^{18}O}. $\alpha_{\mathrm{A/B}} = (1 - 0.0100)/(1 + 0.0020) = 0.98802$.
\end{solution}

\begin{solution}{exo:b3:statistical-thermo-applied:10}
$\Delta\mathrm{ZPE} = \frac12(1 - 1/\sqrt2)(2000 - 3000) = \qty{-146}{cm^{-1}}$; $K = \eu^{146.4/207.2} = 2.0$.
Deuterium goes preferentially to X, the stiffer bond, where its \omterm{def:b3:quantum-model-systems:oscillator}{zero-point
energy} saving is largest.
\end{solution}

\begin{solution}{exo:b3:statistical-thermo-applied:11}
$\Delta_rH^\circ = R\ln10\,(-1.766 + 3.392)/(1/900 - 1/1100) = 19.145 \times 1.626/\num{2.020e-4} =
\qty{154}{kJ/mol}$. The two atoms carry $2 \times \frac52RT$ of enthalpy; the molecule $\frac52RT + RT$
(rotation) plus its vibrational energy $R\theta_{\mathrm V}/(\eu^{\theta_{\mathrm V}/T} - 1)$, slightly less than $RT$:
the difference, about \qty{5}{kJ/mol} at \qty{1000}{K}, adds to $D_0$.
\end{solution}

\begin{solution}{exo:b3:statistical-thermo-applied:12}
Levels 0 ($g = 1$) and $6hcB$ ($g = 5$): $q = 1 + 5\eu^{-x}$, $\langle\varepsilon\rangle/kT = 5x\eu^{-x}/q$, $\langle\varepsilon^2\rangle/(kT)^2 =
5x^2\eu^{-x}/q$, and $C/R = \langle\varepsilon^2\rangle/(kT)^2 - (\langle\varepsilon\rangle/kT)^2 = 5x^2\eu^{-x}/(1 + 5\eu^{-x})^2$. At
\qty{100}{K}, $x = 5.121$: $C/R = 0.783/1.061 = 0.738$.
\end{solution}

\begin{solution}{pb:b3:statistical-thermo-applied:1}
\textbf{1.} $2 \times 2 = 4$: three symmetric (the triplet), one antisymmetric (the singlet).
\textbf{2.} Protons are fermions, so the total wavefunction changes sign on
exchange; the rotational function gives $(-1)^J$: even $J$ goes with the
antisymmetric singlet (para), odd $J$ with the symmetric triplet (ortho).
\textbf{3.} Even levels 1, odd levels 3.
\textbf{4.} At \qty{300}{K} ($T = 3.5\,\theta_{\mathrm R}$) the even and odd rotational sums are nearly equal;
weighted 1 and 3 they give 1:3 (para fraction 0.2507).
\textbf{5.} $F(1) = 2 \times 59.322 - 4 \times 0.0471 = \qty{118.46}{cm^{-1}} = \qty{1.417}{kJ/mol}$.
\textbf{6.} $1/(1 + 9\eu^{-170.4/20.37}) = 1/(1 + 9 \times \num{2.3e-4}) = 0.998$.
\textbf{7.} With only $J = 0$ and 1, half para means $9\eu^{-2\theta_{\mathrm R}/T} = 1$, $T = 2\theta_{\mathrm R}/\ln9 =
0.91\,\theta_{\mathrm R} = \qty{78}{K}$: the competition is between the weight 9 of $J = 1$ and its
Boltzmann factor.
\textbf{8.} 0.75. Without a catalyst the conversion needs days, the liquefaction
hours.
\textbf{9.} $(0.75 - 0.002) \times 1.417 = \qty{1.060}{kJ/mol}$.
\textbf{10.} $1.060/\num{2.016e-3} = \qty{526}{kJ/kg}$.
\textbf{11.} $0.905/\num{2.016e-3} = \qty{449}{kJ/kg}$.
\textbf{12.} The heat of conversion is 1.17 times the enthalpy of vaporisation.
\textbf{13.} It is produced inside the liquid, not brought in through the walls.
\textbf{14.} $x_{\mathrm o}(t) = x_{\mathrm{eq}} + (0.75 - x_{\mathrm{eq}})\eu^{-kt}$, with $x_{\mathrm{eq}} = 0.002$.
\textbf{15.} $0.002 + 0.748\eu^{-0.274} = 0.571$.
\textbf{16.} $(0.75 - 0.571) \times 1.417 = \qty{0.254}{kJ/mol}$.
\textbf{17.} $0.254/0.905 = 0.28$: 28\,\% of the tank in a day.
\textbf{18.} After \qty{168}{h}, $x_{\mathrm o} = 0.112$, heat \qty{0.904}{kJ/mol}, equal to the enthalpy of
vaporisation: in this model the tank is empty. The estimate overstates the
loss because the evaporated molecules carry their unconverted \omterm{def:b3:statistical-thermo-applied:spin-isomers}{ortho hydrogen}
away; the real loss is smaller but still large.
\textbf{19.} $t_{1/2} = \ln2/0.0114 = \qty{61}{h}$.
\textbf{20.} In the liquefier, on catalyst beds at successive temperatures, so that the
refrigerator removes the conversion heat before the liquid is stored.
\textbf{21.} Yes: the equilibrium para fraction at \qty{300}{K} is 0.2507.
\textbf{22.} In normal hydrogen the two forms cannot interconvert, so each one's
rotation freezes on its own ladder of levels; in equilibrium hydrogen part of
the heat supplied converts para into ortho.
\textbf{23.} About 3.57 near \qty{50}{K}: besides exciting rotation, the heat converts
molecules from $J = 0$ to $J = 1$, \qty{118}{cm^{-1}} higher, a reaction with a positive
enthalpy.
\textbf{24.} 1.50: rotation is frozen in both forms (para in $J = 0$, ortho in $J = 1$).
\textbf{25.} \textbf{Heat of conversion / enthalpy of vaporisation $\approx 1.2$ for normal
hydrogen}: its conversion alone could evaporate the whole tank, which is why
hydrogen is converted to para before storage.
\end{solution}
